\documentclass[12pt]{article}

\usepackage{listings}
\usepackage{graphicx}
\usepackage{float}
\usepackage{hyperref}
\usepackage{ifthen}

\usepackage{amssymb} 
\usepackage{amsmath}
\usepackage{amsthm}
\usepackage{sansmath}
\usepackage{epsfig}

\usepackage[T1]{fontenc}
\usepackage{lmodern}
\renewcommand*\familydefault{\sfdefault}

\usepackage[margin=2cm]{geometry}

\usepackage[parfill]{parskip}
\linespread{1.2} 

\usepackage{xcolor}

\lstdefinestyle{code}{  
    commentstyle=\color{gray},
    keywordstyle=\color{blue},
    numberstyle=\tiny\color{gray},
    stringstyle=\color{teal},
    ndkeywordstyle=\color{purple},
    basicstyle=\ttfamily\footnotesize,
    breakatwhitespace=false,         
    breaklines=true,                 
    captionpos=t,                    
    keepspaces=true,                 
    numbers=left,                    
    numbersep=8pt,                  
    showspaces=false,                
    showstringspaces=false,
    showtabs=false,                  
    tabsize=4,
    frame=single,
    breaklines=true
}

\lstset{style=code}

\newcounter{question}
\setcounter{question}{1}
\newcounter{subquest}

\newcommand{\question}[1]{
    \stepcounter{question} 
    \ifnum\value{question}>1 \fi
    \vspace{.5em}
    \textbf{ \Large\thequestion \ - #1} \\
    \vspace{.5em} 
    \hrulefill \\
    \setcounter{subquest}{0}\ }

\newcommand{\subquestion}[1][true]{
    \stepcounter{subquest} 
    \ifthenelse{\equal{#1}{true} \and \value{subquest}>1}{\newpage}{}
    \vspace{1em}
    \textbf{\large Question \thequestion.\thesubquest}
    \vspace{.5em}\ \\}

\newcommand{\solution}
    {\hrulefill \par\vspace{0.5em}\noindent\emph{Solution.}\ }
    {\par\vspace{1em}}

\newcommand{\startappendix}{%
    \newpage
    \vspace{1em}
    {\Large\textbf{Appendix}}

    The complete code used for this assignment is provided in the appendix for reference. Files can be accessed directly at this \href{https://github.com/sjsusu/numerical_analysis}{GitHub repository}.
}

\usepackage{fancyhdr}
\setlength{\headheight}{15pt}
\pagestyle{fancy} 
\lhead{Math 5335: \emph{Numerical Analysis}}
\rhead{Homework 3}
\cfoot{}
\rfoot{\thepage}

\begin{document}
\setcounter{page}{2}
\question{Implementation}
Implement the shooting method with the following specifications:
\begin{itemize}
    \item[(a)] Use RK4 from HW2 as IVP solver.
    \item[(b)] Use secant method as root finder.
    \item[(c)] Report results in table, logging initial guesses $s_0, s_1$ for secant method, all intermediate values of $s_0, s_1, s_2$ for each iteration, the number of secant iterations used, and the final value of $s$ obtained.  
\end{itemize}
Use the following stopping conditions for the secant method:
\[ |F(s_k)|<10^-10 \quad \text{and} \quad |s_k - s_{k-1}| < 10^{-10} \]
Also, add a reasonable maximum iteration cap.

\solution
temp

\hrulefill

\question{Two Independent Verification Checks}
After obtaining $s$ and the numerical solution $y_{N+1}(x)$, verify correctness of your implementation with the following checks:
\begin{itemize}
    \item[(i)] Check the residual $F$ by reporting the final residual magnitude $|F(s)|$.
    \item[(ii)] Check solution error by comparing $y_{N+1}(x)$ to the exact solution $y(x)$. Compute the $L^2$ error:
    \[ e_{L^2} = \sum_i (y_i(x_i) - y(x_i))^2 \]
\end{itemize}

\solution
temp

\hrulefill

\question{Step Size Refinement and Order of Accuracy}
Perform a refinement study using the following step sizes:
\[h=\left\{\frac{\pi}{10}, \frac{\pi}{20}, \frac{\pi}{40}, \frac{\pi}{80}, \frac{\pi}{160}\right\}\]
For each $h$, run the full shooting procedure and compute the $L^2$ error and residual. 

\begin{itemize}
    \item[(a)] Create a table with the following columns:
    \[ h, \quad s, \quad |F(s)|, \quad e_{L^2}(h) \]
    \item[(b)] For $h=\frac{\pi}{160}$, plot $y_i$ and $y(x)$ on the same graph.
    \item[(c)] Make a log-log plot of $e_{L^2}$ vs $h$. Estimate the observed order of accuracy.
\end{itemize}

\solution
temp

\end{document}